(1) 导电不可压黏液在两固定平板 \(y=\pm h\) 间沿 \(x\) 方向稳定层流，黏度 \(\mu>0\)，电导率 \(\sigma>0\)，密度 \(\rho\)。给定恒定压力梯度 \(\partial_xp=-G<0\) 与横向磁感应强度 \(B\hat y\)。磁Reynolds数很小，外部电极短接以维持 \(E_z=0\)，使局部 \(j_z=\sigma Bu\)。从Ohm定律、磁Lorentz力和Stokes动量方程建立 \(u(y)\) 的边值问题。
(2) 引入Hartmann数 \(Ha=Bh\sqrt{\sigma/\mu}\)，无滑移 \(u(\pm h)=0\)，求速度分布、每单位横向宽度的总体积流率 \(Q\) 和平均速度，证明 \(Ha\to0\) 恢复经典平板Poiseuille解。
(3) 对 \(Ha\gg1\) 求芯部速度、两侧边界层厚度 \(\delta\)、流量首阶渐近；说明为何大磁场提高流动电阻并检验 \(j_z\) 的流向及磁阻力方向。
(4) 将定常流动方程乘以速度在截面上积分，从而严格证明外部压强功等于黏性剪切耗散加Ohm焦耳耗散，显式给出两个积分的量纲；讨论若电路开路导致 \(E_z\) 自行建立，解为何发生变化。
